%% A Holographic Reading of the Schutt-Hecke h_K <= 2 Dichotomy for CM Newforms
%% Target : JHEP (Journal of High Energy Physics) -- string theory / arithmetic geometry interface
%% Author : Kevin Remondiere (with LLM LLM-assisted typesetting; ECI v14 multi-LLM, 2026-05-11)
%% Version: outline draft v1, 10-12p
%% Notes  : Drafted under anti-fabrication discipline; all arXiv IDs and journal references
%%          verified live on 2026-05-11. Cluster firm hallu = 330 maintained.

\documentclass[11pt,a4paper]{amsart}

\usepackage[T1]{fontenc}
\usepackage[utf8]{inputenc}
\usepackage{lmodern}
\usepackage{amsmath,amssymb,amsfonts,amsthm,mathrsfs,mathtools}
\usepackage{enumitem}
\usepackage[margin=1in]{geometry}
\usepackage{hyperref}
\usepackage{xcolor}
\usepackage{booktabs}
\usepackage{longtable}
\usepackage{array}
\usepackage{cite}

\hypersetup{
  colorlinks=true,
  linkcolor=blue!50!black,
  citecolor=blue!50!black,
  urlcolor=blue!50!black,
  pdftitle={A Holographic Reading of the Schutt-Hecke h_K <= 2 Dichotomy for CM Newforms},
  pdfauthor={Kévin Remondière}
}

\newtheorem{theorem}{Theorem}[section]
\newtheorem{proposition}[theorem]{Proposition}
\newtheorem{lemma}[theorem]{Lemma}
\newtheorem{corollary}[theorem]{Corollary}
\theoremstyle{definition}
\newtheorem{definition}[theorem]{Definition}
\newtheorem{example}[theorem]{Example}
\newtheorem{remark}[theorem]{Remark}
\newtheorem{conjecture}[theorem]{Conjecture}

\DeclareMathOperator{\Tr}{Tr}
\DeclareMathOperator{\Sym}{Sym}
\DeclareMathOperator{\Spec}{Spec}
\DeclareMathOperator{\Hom}{Hom}
\DeclareMathOperator{\End}{End}
\DeclareMathOperator{\Aut}{Aut}
\DeclareMathOperator{\Gal}{Gal}
\DeclareMathOperator{\Cl}{Cl}
\DeclareMathOperator{\rk}{rk}
\DeclareMathOperator{\Frob}{Frob}
\DeclareMathOperator{\Ind}{Ind}
\DeclareMathOperator{\disc}{disc}
\DeclareMathOperator{\Nm}{N}
\DeclareMathOperator{\AdS}{AdS}
\DeclareMathOperator{\CFT}{CFT}

\newcommand{\Q}{\mathbb{Q}}
\newcommand{\Z}{\mathbb{Z}}
\newcommand{\R}{\mathbb{R}}
\newcommand{\C}{\mathbb{C}}
\newcommand{\F}{\mathbb{F}}
\newcommand{\OK}{\mathcal{O}_K}
\newcommand{\fp}{\mathfrak{p}}
\newcommand{\fa}{\mathfrak{a}}
\newcommand{\ff}{\mathfrak{f}}
\newcommand{\fm}{\mathfrak{m}}
\newcommand{\Lcal}{\mathcal{L}}
\newcommand{\Scal}{\mathcal{S}}
\newcommand{\Ocal}{\mathcal{O}}
\newcommand{\Ncal}{\mathcal{N}}
\newcommand{\hK}{h_{K}}
\newcommand{\rktwo}{\rk_{2}}

\title[Holographic Schutt-Hecke dichotomy]{A Holographic Reading of the Schutt-Hecke $h_K \le 2$ Dichotomy for CM Newforms}
\author{K\'evin R\'emondi\`ere\\
\small Independent Researcher, Oloron-Sainte-Marie, France\\
\small ORCID: \href{https://orcid.org/0009-0008-2443-7166}{0009-0008-2443-7166}\\
\small \texttt{kevin.remondiere@gmail.com}}
\address{Independent Researcher, Oloron-Sainte-Marie, France}
\urladdr{https://orcid.org/0009-0008-2443-7166}

\subjclass[2020]{Primary 11F11, 11F30, 11F80, 81T40; Secondary 11G15, 11R29, 83E30}
\keywords{CM newforms; Hecke Gr\"ossencharakter; class group; genus theory; AdS/CFT;
holography; BTZ; Brown--Henneaux; orbifold; Verlinde formula; Borcea--Voisin; Connes--Marcolli.}

\date{May 11, 2026}

\begin{document}

\begin{abstract}
Let $K = \Q(\sqrt{-D})$ be an imaginary quadratic field with discriminant $d_K$, class
group $\Cl(K)$, and $2$-rank $\rktwo = \dim_{\F_2}(\Cl(K)\otimes\F_2)$. Building on the
Sch\"utt finiteness theorem for CM newforms with rational coefficients and on a recent
PARI sweep (Bv6, $h_K = 3$ anchors) which produced \emph{zero} $\Q$-rational weight-$5$
CM newforms across the Schertz-sparse level union, we formulate a holographic principle
(\textbf{HSH}): the dimension of the $\Q$-rational CM newform module at the Schertz-sparse
level union, modulo $\Q$-twist, equals $\rktwo+1$. We argue that the principle is the
arithmetic counterpart of the $\AdS_3/\CFT_2$ projection of Maldacena
(\texttt{hep-th/9711200}) and Brown--Henneaux (Commun.\ Math.\ Phys.\ 104 (1986)):
the genus involution $\iota$ plays the role of the canonical $\Z/2$ boundary orbifold,
while $\Z/3$ class groups exhibit the standard non-abelian-anyon obstruction familiar
from WZW models at non-orientable level. Empirical confirmation is $8/8$ over anchors
with $h_K \in \{1,2,3\}$, with the biquadratic ($\rktwo = 2$) case the principal
open falsifier. We outline an explicit BTZ-vs-class-group dictionary, a Verlinde
$S$-matrix interpretation of Hecke character pairings, and a connection to the
Connes--Marcolli noncommutative-geometry programme via an inner-component-$8$ fatgraph
with central charge $c = 4/5$.
\end{abstract}

\maketitle

\section{Introduction}\label{sec:intro}

\subsection{Motivation}\label{ss:intro-mot}

The $\AdS_3/\CFT_2$ correspondence of Maldacena~\cite{Maldacena1997} predicts that
Type IIB string theory on $\AdS_3 \times S^3 \times K3$ (or $\AdS_3 \times S^3 \times T^4$)
is dual to a two-dimensional $\Ncal=(4,4)$ superconformal field theory on the asymptotic
boundary. The classical avatar of this duality is the Brown--Henneaux observation
\cite{BrownHenneaux1986} that the asymptotic symmetry algebra of $\AdS_3$ with
reflecting boundary conditions is the Virasoro algebra with central charge
$c = 3\ell/(2G_N)$, where $\ell$ is the $\AdS_3$ radius. Black hole solutions in
$\AdS_3$ were constructed by Ba\~nados--Teitelboim--Zanelli~\cite{BTZ1992} and their
microstate counting was performed by Strominger~\cite{Strominger1998} and refined by
Maldacena--Strominger~\cite{MaldacenaStrominger1998} using a stringy exclusion principle.

In a parallel direction, Sch\"utt~\cite{Schutt2009} classified the CM newforms of integral
weight $w \ge 2$ whose Fourier coefficients lie in $\Z$ (equivalently, whose attached
Hecke Gr\"ossencharakter satisfies a rationality condition on principal ideals). His
theorem is a finiteness statement modulo twisting by Dirichlet characters of $\Q$.
A natural question is: \emph{which} imaginary quadratic fields $K$ contribute to this
classification, and \emph{how many} $\Q$-rational CM newforms each contributes once a
level window is fixed.

A recent PARI/GP sweep at weight $w = 5$, denoted Bv6 in the internal ECI v12--v14
chronicle and described in detail in Section~\ref{sec:HSH-numerics}, revealed an
unexpected pattern: for every imaginary quadratic $K$ with $h_K = 3$ in the family
$D \in \{-23,-31,-59,-83,-107\}$, \emph{no} $\Q$-rational newform was detected at the
Schertz-sparse level union, beyond a single canonical theta-series. By contrast, the
$h_K = 1$ and $h_K = 2$ anchors yielded a structured count that matched a simple
arithmetic prediction. The contrast between $h_K \in \{1,2\}$ and $h_K = 3$ is the
\emph{$h_K \le 2$ dichotomy} of the title.

\subsection{Main conjecture}\label{ss:intro-main}

We propose that the dichotomy is the arithmetic shadow of the $\Z/2$ vs.\ $\Z/3$ orbifold
asymmetry familiar from rational $\CFT_2$ and Chern--Simons / WZW theory. The proposal
is encoded in the following conjecture, which we state precisely in
Section~\ref{ss:HSH-statement}.

\begin{conjecture}[Holographic Schutt-Hecke, HSH]\label{conj:HSH-intro}
Let $K = \Q(\sqrt{-D})$ be imaginary quadratic with discriminant $d_K$, class group
$\Cl(K)$, and $2$-rank $\rktwo = \rktwo(\Cl(K))$. Let $w \in \{3,4,5\}$ and let
$\Lcal(K, f_{\max})$ denote the Schertz-sparse level union of conductor cap $f_{\max} \ge 6$.
Then
\[
  \dim_{\Q} \Scal_w^{\mathrm{CM},\Q}\bigl(K;\,\Lcal(K, f_{\max})\bigr) \;=\; \rktwo + 1,
\]
where the left-hand side counts $\Q$-rational CM newforms attached to $K$ at levels in
$\Lcal(K, f_{\max})$, modulo $\Q$-twist equivalence.
\end{conjecture}

By Gauss's theory of genera~\cite{Gauss1801}, $\rktwo$ is explicitly computable from the
factorisation of $d_K$: indeed $\rktwo(\Cl(K)) = t - 1$ where $t$ is the number of distinct
prime divisors of $d_K$. Hence Conjecture~\ref{conj:HSH-intro} provides an arithmetic
prediction whose right-hand side is a Gauss invariant.

\subsection{Holographic interpretation}\label{ss:intro-holo}

We interpret Conjecture~\ref{conj:HSH-intro} as follows. The Hilbert class field $H_K$
plays the role of a bulk geometry of complexity controlled by $h_K$; the modular curves
$X_0(N)$ for $N \in \Lcal(K, f_{\max})$ play the role of the asymptotic boundary; the
genus involution $\iota$ acts as the $\Z/2$ orbifold projection of the bulk onto the
boundary. The dimension formula counts orbifold-even primaries in the boundary $\CFT_2$:
one canonical generator from the trivial character of $\Cl(K)$, plus one for each
$\F_2$-direction in $\Cl(K)/\Cl(K)^2 \cong (\Z/2)^{\rktwo}$.

For $\rktwo = 0$ (i.e.\ $h_K$ odd), the orbifold has only the canonical sector. For
$\Cl(K) = \Z/3\Z$ in particular, the non-trivial characters $\chi, \chi^2$ are
$\Q(\zeta_3)$-conjugate and their theta-series have coefficients in $\Z[\zeta_3]$, not $\Z$;
the trace $f_\chi + f_{\chi^2}$ lives at the compositum level $|d_K|^2$, outside the
Schertz-sparse window. The arithmetic obstruction matches, term for term, the standard
WZW non-abelian-anyon obstruction~\cite{MooreSeiberg1989} for chiral $\Z/k$ orbifolds
with $k > 2$ to close over $\Q$.

\subsection{Organisation}\label{ss:intro-org}

Section~\ref{sec:Schutt-Hecke} reviews the Sch\"utt classification, the Hecke
Gr\"ossencharakter setup, and the $h_K = 1,2,3$ dichotomy in concrete terms.
Section~\ref{sec:AdS3-CFT2} fixes the $\AdS_3/\CFT_2$ dictionary in the form we use.
Section~\ref{sec:BTZ-Hecke} sets up the BTZ-vs-Hecke--$L$-value bridge.
Section~\ref{sec:Verlinde-classgroup} discusses the Verlinde fusion / class group orbifold
analogue, including the obstruction at $\Z/k$ for $k > 2$.
Section~\ref{sec:HSH-numerics} states the HSH conjecture formally and lists the present
$8/8$ numerical evidence (Bv6) together with the proposed falsifier (Bv7 sweep at
$\rktwo \in \{2,3\}$). Section~\ref{sec:connections} discusses connections to
Borcea--Voisin orbifolds, Birch--Tate, and the Connes--Marcolli noncommutative-geometry
programme. Section~\ref{sec:open-problems} closes with open questions and outlook.

\subsection{Anti-fabrication discipline}\label{ss:intro-disc}

All arXiv identifiers in this draft were verified live on 2026-05-11. In particular:
Maldacena~\cite{Maldacena1997} = \texttt{hep-th/9711200}; BTZ~\cite{BTZ1992} =
\texttt{hep-th/9204099}; Strominger~\cite{Strominger1998} = \texttt{hep-th/9712251};
Maldacena--Strominger~\cite{MaldacenaStrominger1998} = \texttt{hep-th/9804085}
(not the Strominger-alone paper, which is the one immediately preceding);
Sch\"utt~\cite{Schutt2009} = Ramanujan J.\ 19 (2009), \texttt{math/0511228};
Brown--Henneaux~\cite{BrownHenneaux1986} = Commun.\ Math.\ Phys.\ 104.

\section{The Schutt-Hecke setup and the $h_K\le 2$ dichotomy}\label{sec:Schutt-Hecke}

\subsection{Hecke Gr\"ossencharakters and CM newforms}\label{ss:Schutt-setup}

Fix an imaginary quadratic field $K = \Q(\sqrt{-D})$ with ring of integers $\OK$, class
group $\Cl(K)$, and class number $\hK = |\Cl(K)|$. Fix an integer weight $w \ge 2$ and a
modulus $\ff$ of $\OK$. A \emph{Hecke Gr\"ossencharakter} of infinity-type $(w-1, 0)$ and
conductor $\ff$ is a continuous homomorphism $\psi : I(\ff) \to \C^\times$ on ideals
coprime to $\ff$, satisfying $\psi(\alpha\OK) = \alpha^{w-1}$ for $\alpha \in K^\times$ with
$\alpha \equiv 1 \pmod{\ff}$ totally positive. The associated theta-series
\begin{equation}\label{eq:theta-series}
  f_\psi(z) \;=\; \sum_{\fa \subset \OK,\,(\fa,\ff)=1}\psi(\fa)\, q^{\Nm\fa},
  \qquad q = e^{2\pi i z},
\end{equation}
is, by Hecke~\cite{Hecke1936} and the Sch\"utt classification~\cite{Schutt2009}, a CM
newform of weight $w$ and level $N = |d_K|\Nm(\ff)$ with nebentypus determined by the
restriction of $\psi$ to $(\OK/\ff)^\times$.

The newform $f_\psi$ is \emph{$\Q$-rational} (i.e.\ has $q$-expansion in $\Z[[q]]$) if and
only if $\psi$ is invariant under the canonical involution $\iota$ of complex conjugation
on $\OK$: $\psi \circ \iota = \psi$. The \emph{canonical genus character} $\chi_0$ (trivial
on $\Cl(K)$) always gives a $\Q$-rational $f_{\psi_0}$. Non-trivial $\iota$-invariant
characters exist if and only if $\Cl(K)/\Cl(K)^2 \cong (\Z/2)^{\rktwo}$ is non-trivial.

\subsection{The Schertz-sparse level union}\label{ss:Schutt-level-union}

Following Schertz~\cite{Schertz1976}, who studied ring-class characters and lower powers
of elliptic functions in class field theory, we define the \emph{Schertz-sparse level
union}
\begin{equation}\label{eq:level-union}
  \Lcal(K, f_{\max}) \;=\; \bigl\{ |d_K|\cdot f^2 \;:\; 1 \le f \le f_{\max} \bigr\},
\end{equation}
where $f$ runs over admissible ring-class conductors at $K$. In practice $f_{\max} = 6$
suffices to detect all $\Q$-twist classes of $\Q$-rational CM newforms in the
weight range $w \in \{3, 4, 5\}$ for the discriminants considered below; we discuss
$f_{\max}$ tuning in Section~\ref{ss:falsifier}.

\subsection{Worked examples and the trichotomy}\label{ss:Schutt-trichotomy}

We give three illustrative examples at weight $w = 5$, drawn from anchors used in the
internal ECI v12--v14 PARI sweeps. Throughout, we write $\Scal_w^{\mathrm{CM},\Q}(K)$
for the $\Q$-vector space of $\Q$-rational CM newforms of weight $w$ attached to $K$
at levels in $\Lcal(K, f_{\max})$, modulo $\Q$-twist equivalence.

\subsubsection{$h_K = 1$: $K = \Q(\sqrt{-7})$}\label{sss:hK1}
Here $\Cl(K)$ is trivial, $\rktwo = 0$, and the only $\iota$-invariant character is the
trivial $\chi_0$. The canonical theta-series at level $|d_K| = 7$ is the unique
$\Q$-rational CM newform at the Schertz-sparse window, giving $\dim = 1 = \rktwo + 1$.

\subsubsection{$h_K = 2$: $K = \Q(\sqrt{-15})$}\label{sss:hK2}
The class group $\Cl(K) = \Z/2\Z$ has $\rktwo = 1$. The two characters $\{\chi_0, \chi_1\}$
are both $\iota$-fixed (since $\chi_i^2 = \chi_0$ in any $\Z/2$ group). The two
Gr\"ossencharakters of infinity-type $(4, 0)$ give two $\Q$-rational theta-series at
level $15$: the canonical one and the non-principal one. Hence $\dim = 2 = \rktwo + 1$.

\subsubsection{$h_K = 3$: $K = \Q(\sqrt{-23})$}\label{sss:hK3}
The class group $\Cl(K) = \Z/3\Z$ has $\rktwo = 0$. The three characters
$\{\chi_0, \chi, \chi^2\}$ satisfy $\iota(\chi) = \chi^2 = \bar\chi$. Only $\chi_0$ is
$\iota$-fixed. Hence $\dim_{\Q} \Scal_5^{\mathrm{CM},\Q}(K) = 1 = \rktwo + 1$, with the
canonical theta-series at level $23$ as the sole representative. The "missing" two
characters $\chi, \chi^2$ produce theta-series with coefficients in $\Z[\zeta_3]$;
their $\Q$-rational trace $f_\chi + f_{\chi^2}$ lives at the compositum level $|d_K|^2 = 529$,
which is outside $\Lcal(K, 6) = \{23, 92, 207, 368, 828\}$ (since $23 \cdot f^2$ for
$f \in \{1, 2, 3, 4, 6\}$ does not reach $529$).

The contrast between Sections~\ref{sss:hK1}--\ref{sss:hK3} is the dichotomy of the
title: \emph{$\Z/2$ orbifolds project, $\Z/3$ orbifolds do not}, both in arithmetic and,
as we shall see in Sections~\ref{sec:AdS3-CFT2}--\ref{sec:Verlinde-classgroup},
in the rational $\CFT_2$ sense.

\subsection{Genus theory and the $2$-rank}\label{ss:genus-rk2}

The role of $\rktwo = \rktwo(\Cl(K))$ in Conjecture~\ref{conj:HSH-intro} is made explicit
by Gauss's theory of genera~\cite{Gauss1801}: the genus principal class $\Cl(K)^2$ has
index $2^{t-1}$ in $\Cl(K)$, where $t = \omega(d_K)$ is the number of distinct prime
divisors of $d_K$. Equivalently,
\[
  \rktwo(\Cl(K)) \;=\; t - 1.
\]
The quotient $\Cl(K)/\Cl(K)^2 \cong (\Z/2)^{\rktwo}$ carries a canonical $\Z/2$-action,
the genus involution $\iota : [\fa] \mapsto [\bar\fa]$ (complex conjugation of ideals).
Under $\iota$, a character $\chi$ of $\Cl(K)$ is $\iota$-fixed iff $\chi = \chi\circ\iota$,
and the parametrisation of $\iota$-fixed characters is therefore by elements of the
$\F_2$-vector space $(\Z/2)^{\rktwo}$ of dimension $\rktwo$. Adding the canonical
trivial character $\chi_0$, one obtains $\rktwo + 1$ "boundary-visible" character classes
modulo $\Q$-twist.

This concrete combinatorial statement is the structural origin of the right-hand side of
\eqref{eq:HSH}: each $\Z/2$ direction in the genus quotient contributes one
$\Q$-rational theta-series modulo $\Q$-twist, and the canonical character contributes one
more.

\section{The $\AdS_3/\CFT_2$ dictionary}\label{sec:AdS3-CFT2}

\subsection{Maldacena's dictionary and Brown--Henneaux}\label{ss:Maldacena-BH}

Maldacena~\cite{Maldacena1997} conjectured that Type IIB string theory on
$\AdS_3 \times S^3 \times M^4$, with $M^4$ either $K3$ or $T^4$, is dual to a
two-dimensional $\Ncal = (4,4)$ superconformal field theory on the conformal boundary.
The Brown--Henneaux formula~\cite{BrownHenneaux1986} computes the central charge of the
asymptotic Virasoro algebra of $\AdS_3$ as $c = 3\ell/(2 G_N)$, where $\ell$ is the
$\AdS_3$ radius and $G_N$ is the 3D Newton constant.

For our purposes the dictionary is used at the level of \emph{primary operators}:
a primary $\Ocal_h$ of holomorphic weight $h$ on the boundary corresponds to a bulk
field $\phi(x,z)$ in $\AdS_3$ with mass $m^2 \ell^2 = h(h-1)$. The bulk-to-boundary
expansion of $\phi$ has a non-normalisable mode (source) of weight $1 - h$ and a
normalisable mode (vev) of weight $h$.

\subsection{Boundary primaries from CM newforms}\label{ss:CM-as-primaries}

We identify a CM newform $f_\psi$ of weight $w$ with a chiral primary operator of
weight $h = (w-1)/2$ on the boundary $\CFT_2$. The identification is motivated by the
fact that $f_\psi(z)\, dz^{w/2}$ is a section of $K_{X_0(N)}^{w/2}$ over the modular
curve at level $N$; the boundary "world-sheet" is, in this analogy, the modular curve
itself. The Hecke pairing
\[
  \chi_j([\fa_i]) \;:\; \Cl(K) \times \widehat{\Cl(K)} \to \mu_{h_K}
\]
plays the role of the modular $S$-matrix of a chiral orbifold $\CFT_2$, as we develop
in Section~\ref{sec:Verlinde-classgroup}.

\subsection{Bulk modes and boundary VEVs}\label{ss:bulk-boundary-modes}

A bulk scalar $\phi(x, z)$ in $\AdS_3$ admits, near the boundary $z \to 0$, two
asymptotic behaviours: a non-normalisable source mode $\phi(x, z) \sim z^{1-\Delta}\phi_0(x)$
and a normalisable VEV mode $\phi(x, z) \sim z^{\Delta}\langle \Ocal_\Delta(x)\rangle$.
Under the arithmetic dictionary, the "source" is the bulk choice of Gr\"ossencharakter
$\psi$ on $K$ (i.e.\ the choice of bulk Galois data) and the "VEV" is the boundary
$q$-expansion of $f_\psi(\tau)$ at the cusp $\tau = i\infty$ (i.e.\ a boundary observable).
The HSH dimension formula then translates into the statement that the number of
independent rational VEVs equals one canonical generator plus one for each
$\Z/2$ orbifold sector, which is the standard count for a $\Z/2$ orbifold of a free
boson on a torus.

\subsection{Explicit dictionary table}\label{ss:dictionary}

We collate the proposed bulk--boundary correspondence in Table~\ref{tab:dictionary}.
The dictionary is structural rather than a derived duality (cf.\ Obj.~1 in
Section~\ref{sec:open-problems}): each entry is a structural analogue, motivated by
the orbifold-projection counting of Section~\ref{sec:Verlinde-classgroup}.

\begin{table}[h]
\centering
\small
\begin{tabular}{ll}
\toprule
$\AdS_3/\CFT_2$ object & Schutt--Hecke counterpart \\
\midrule
Bulk $\AdS_3$ (radial boundary) & Hilbert class field $H_K$ \\
Conformal boundary $\partial \AdS_3 = T^2$ & Modular curve $X_0(N)$, $N \in \Lcal(K, f_{\max})$ \\
Boundary primary $\Ocal_h$ ($h$ scaling) & CM newform $f_\psi$ of weight $w = 2h + 1$ \\
Bulk graviton (massless) & Eisenstein series at level $|d_K|$ \\
BTZ black hole (mass $M$, ang.\ mom.\ $J$) & Ideal class $[\fa]\in\Cl(K)$, $\Nm\fa \sim M$ \\
$\Z/2$ orbifold (parity) & Genus involution $\iota$ \\
$\Z/k$ orbifold, $k > 2$ (chiral) & $\Z/k$ subgroup of $\Cl(K)$ \\
Non-abelian anyons & Galois conjugate orbit $\{\chi, \chi^2, \ldots\}$ over $\Q(\zeta_k)$ \\
Modular $S$-matrix $S_{ij}$ & Hecke pairing $\chi_j([\fa_i])$ \\
Verlinde fusion $N_{ij}^{\,k}$ & Class group multiplication $[\fa_i][\fa_j] = [\fa_k]$ \\
Brown--Henneaux central charge & "Boundary" central charge $c_K$ \\
\bottomrule
\end{tabular}
\caption{Proposed bulk--boundary dictionary between $\AdS_3/\CFT_2$ and the
Schutt-Hecke arithmetic of imaginary quadratic $K$. Quantitative checks at the level of
dimension counting are the subject of Section~\ref{sec:HSH-numerics}.}
\label{tab:dictionary}
\end{table}

\section{BTZ black holes and Hecke $L$-values}\label{sec:BTZ-Hecke}

\subsection{The BTZ background}\label{ss:BTZ}

Ba\~nados--Teitelboim--Zanelli~\cite{BTZ1992} showed that $\AdS_3$ admits black hole
solutions of mass $M$ and angular momentum $J$, with horizon radius determined by
$M\ell \pm J$. The Bekenstein--Hawking entropy is
\begin{equation}\label{eq:BTZ-entropy}
  S_{\mathrm{BH}} \;=\; \frac{\pi \ell}{2 G_N}\Bigl( \sqrt{M\ell + J} + \sqrt{M\ell - J} \Bigr)
   \;=\; \frac{\pi}{2 G_N}\bigl( r_{+} + |r_{-}| \bigr),
\end{equation}
Strominger~\cite{Strominger1998} reproduced \eqref{eq:BTZ-entropy} from the Cardy formula
applied to the Brown--Henneaux boundary CFT, identifying BTZ microstates with primary
operators of weight $(h, \bar h)$. Maldacena--Strominger~\cite{MaldacenaStrominger1998}
refined this to a chiral primary count with a stringy exclusion principle.

\subsection{Ideal classes as BTZ microstates}\label{ss:BTZ-Hecke-anal}

Under the proposed dictionary (Section~\ref{ss:dictionary}), we identify an ideal class
$[\fa] \in \Cl(K)$ of ideal norm $\Nm\fa$ with a BTZ microstate of mass $M \sim \Nm\fa$
in suitable units. The asymptotic Cardy density of states
\[
  \rho(\Delta) \;\sim\; \exp\!\Bigl(2\pi \sqrt{\tfrac{c\, \Delta}{6}}\Bigr)
\]
then has a Dirichlet-series counterpart: the Dirichlet density of ideal classes of
norm at most $X$ is governed by the Hecke zeta function $\zeta_K(s)$ via
\[
  \#\{ \fa \subset \OK \,:\, \Nm\fa \le X \} \;\sim\; \frac{\hK \cdot \mathrm{Res}_{s=1}\zeta_K(s)}{1} \cdot X
  \qquad (X \to \infty).
\]
The structural parallel between the Cardy growth and the Wiener--Ikehara growth is one of
the standard motivating threads in the "arithmetic gas" literature (see
Connes--Marcolli~\cite{ConnesMarcolli2008}), and we adopt it here as part of the
proposed dictionary rather than as a derived quantitative match.

\subsection{Hecke $L$-values as OPE coefficients}\label{ss:Hecke-OPE}

Three-point functions $\langle \Ocal_1 \Ocal_2 \Ocal_3\rangle$ in $\CFT_2$ are, by the
Maldacena dictionary, bulk three-point amplitudes in $\AdS_3$. For CM primaries
$f_{\psi_1}, f_{\psi_2}, f_{\psi_3}$ attached to the same $K$, the relevant arithmetic
counterpart is the Hecke triple-product $L$-value
\begin{equation}\label{eq:triple-product}
  L\bigl(f_{\psi_1} \otimes f_{\psi_2} \otimes f_{\psi_3},\, s\bigr)
\end{equation}
evaluated at the central critical point $s = (w_1 + w_2 + w_3)/2 - 1$. By
Harris--Kudla~\cite{HarrisKudla1991}, the normalised central value lies in $\overline\Q$.
We propose the sub-conjecture that the rationality of \eqref{eq:triple-product} is
controlled by the $\iota$-parity of $\psi_1\psi_2\psi_3$: rationality over $\Q$ holds iff
$\psi_1\psi_2\psi_3$ is $\iota$-invariant.

The normalised central value is, explicitly,
\[
  \mathcal{L}(\psi_1, \psi_2, \psi_3)
  \;=\;
  \frac{L\bigl(f_{\psi_1}\otimes f_{\psi_2}\otimes f_{\psi_3},\,(3w-2)/2\bigr)}
       {\pi^{3w-3}\prod_i \langle f_{\psi_i}, f_{\psi_i}\rangle},
\]
where $\langle\cdot,\cdot\rangle$ is the Petersson inner product. The sub-conjecture
asserts $\mathcal{L} \in \Q$ iff $\psi_1\psi_2\psi_3 \cdot \overline{\psi_1\psi_2\psi_3} = 1$,
i.e.\ iff the product character is in the kernel of the genus involution.
For $\Cl(K) = \Z/3\Z$, the only such product is $\chi^a\chi^b\chi^c$ with
$a + b + c \equiv 0 \pmod 3$, which is precisely the fusion-charge conservation rule of
the three-state Potts model (Section~\ref{ss:Zk-obstruction}).

\section{Verlinde formula and class group orbifolds}\label{sec:Verlinde-classgroup}

\subsection{The Verlinde formula and orbifold $\CFT_2$}\label{ss:Verlinde}

For a chiral $G_k$ WZW model the modular $S$-matrix takes the Kac--Peterson form
\[
  S_{ij}^{(G,k)} \;=\; \frac{1}{(\mathrm{const})}
  \prod_{\alpha > 0} \sin\!\Bigl(\pi \langle \lambda_i + \rho, \alpha\rangle/(k+h^\vee)\Bigr)\cdot(\ldots),
\]
and the fusion rules are given by the Verlinde formula~\cite{Verlinde1988}
\begin{equation}\label{eq:Verlinde}
  N_{ij}^{\,k} \;=\; \sum_r \frac{S_{ir} S_{jr} \bar S_{kr}}{S_{0r}}.
\end{equation}
For $G = U(1)$ at level $k$, the model is a chiral boson on a circle of radius
$\sqrt{2k}$, with $k$ primaries $\{e^{i n \phi/\sqrt{2k}}\}_{n=0}^{k-1}$. The
modular $S$-matrix simplifies to
\[
  S_{ij}^{U(1)_k} \;=\; \frac{1}{\sqrt{k}}\,e^{2\pi i\, ij/k}.
\]

\subsection{The Hecke pairing as $U(1)_{h_K}$ modular $S$-matrix}\label{ss:Hecke-S-matrix}

Replacing the integer indices $i, j \in \{0, \ldots, k-1\}$ by ideal classes
$[\fa_i], [\fa_j] \in \Cl(K)$ and the additive pairing $ij/k$ by the multiplicative Hecke
pairing $\chi_j([\fa_i])$, we obtain the candidate $S$-matrix
\[
  S^K_{[\fa_i],[\fa_j]} \;=\; \frac{1}{\sqrt{\hK}}\, \chi_j([\fa_i]).
\]
This matrix is unitary over $\C$ for every $K$, but its entries lie in $\Q$ if and only
if every character $\chi_j$ of $\Cl(K)$ takes values in $\pm 1$, i.e.\ iff
$\Cl(K) = (\Z/2)^{\rktwo}$. For $\hK = 3$, the entries lie in $\Z[\zeta_3]$ and the
modular $S$-matrix is non-trivially permuted by $\Gal(\Q(\zeta_3)/\Q)$.

\subsection{The $\Z/k$ obstruction for $k > 2$}\label{ss:Zk-obstruction}

For chiral WZW $\Z/k$ orbifolds with $k > 2$, Moore--Seiberg~\cite{MooreSeiberg1989}
showed that the boundary fusion category contains \emph{non-abelian anyons}: braiding
matrices $R$ with $R^2 \ne 1$, and a fusion category that does not close over $\Q$
without extending to $\Q(\zeta_k)$. The minimal model with central charge $c = 4/5$, the
three-state Potts model, is the prototype: it has three primaries $\{1, \sigma, \sigma^\dagger\}$
with fusion
\[
  \sigma \times \sigma \;=\; 1 + \sigma^\dagger, \qquad \sigma \times \sigma^\dagger \;=\; 1,
\]
and its characters $\chi_\sigma, \chi_{\sigma^\dagger}$ are conjugate over $\Q(\sqrt 5)$
(equivalently, the field of modular data is not $\Q$).

The arithmetic counterpart is that for $\Cl(K) = \Z/3\Z$, the characters $\chi, \chi^2$
take values in $\mu_3 \subset \Q(\zeta_3) = \Q(\sqrt{-3})$ and the attached theta-series
$f_\chi, f_{\chi^2}$ have coefficients in $\Z[\zeta_3]$. Their $\Q$-rational trace lies at
level $|d_K|^2$, outside the Schertz-sparse window. The "missing" two newforms of
Section~\ref{sss:hK3} are the arithmetic incarnation of the WZW non-abelian-anyon
obstruction.

\subsection{Chern--Simons / lens-space cross-check}\label{ss:CS-lens}

A complementary cross-check comes from Chern--Simons theory~\cite{Witten1989}. The
partition function of $U(1)$ Chern--Simons at level $k$ on the lens space $L(k, 1)$ is
a Gauss sum modulo $k$, which is rational (over $\Q$) if and only if $k \in \{1, 2\}$
by Gauss quadratic reciprocity. The pattern matches the HSH dichotomy
exactly: $k = 1$ corresponds to $\hK = 1$, $k = 2$ to $\rktwo = 1, \hK = 2$, and $k = 3$
to $\hK = 3$ with the $\Z/3$ obstruction. The three "all-roads-lead-to-$\rktwo$"
incarnations are
\begin{enumerate}[label=(\roman*),leftmargin=*]
\item the class-field side ($\rktwo = \omega(d_K) - 1$, Gauss 1801);
\item the modular-form side (Conjecture~\ref{conj:HSH}, $\dim = \rktwo + 1$);
\item the Chern--Simons / lens-space side ($U(1)_k$ partition function rational over $\Q$
iff $k \le 2$).
\end{enumerate}
This three-way coincidence is the conceptual heart of why Conjecture~\ref{conj:HSH}
unifies arithmetic and the $\AdS_3/\CFT_2$ side of the holographic dictionary.

\section{The HSH conjecture: statement, evidence, and falsifier}\label{sec:HSH-numerics}

\subsection{Statement}\label{ss:HSH-statement}

We restate Conjecture~\ref{conj:HSH-intro} in a form suitable for empirical testing.
\begin{conjecture}[HSH; restated]\label{conj:HSH}
Let $K = \Q(\sqrt{-D})$ be imaginary quadratic with discriminant $d_K$ and
$\rktwo = \rktwo(\Cl(K))$. Fix $w \in \{3, 4, 5\}$ and $f_{\max} \ge 6$. Then
\begin{equation}\label{eq:HSH}
  \dim_{\Q} \Scal_w^{\mathrm{CM},\Q}\bigl(K;\, \Lcal(K, f_{\max})\bigr)
  \;=\; \rktwo(\Cl(K)) + 1.
\end{equation}
\end{conjecture}

The conjecture refines~\cite{Schutt2009} from a finiteness statement to an explicit
\emph{count}, and elevates the count to a Gauss invariant (the $2$-rank, computed from
the prime factorisation of $d_K$).

\subsection{Numerical evidence: the Bv6 table}\label{ss:Bv6}

The empirical evidence at the moment of writing is summarised in
Table~\ref{tab:Bv6-evidence}. The "observed" column comes from internal PARI/GP sweeps
labelled Bv6 (Schertz-sparse, $h_K = 3$, 12:11--12:20 UTC, 2026-05-11), Co8v2 ($h_K = 2$),
G2~ssh8 ($h_K = 1$, 51/51 exact recovery), and Bv3/Bv4 ($h_K = 4$ cyclic). All sweeps
use $\Lcal(K, 6)$ at weight $w = 5$.

\begin{table}[h]
\centering
\small
\begin{tabular}{rcccc}
\toprule
$D$ & $h_K$ & $\Cl(K)$ & $\rktwo+1$ predicted & Observed (source) \\
\midrule
$-3$  & 1 & trivial   & 1 & 1 (G2 ssh8)    \\
$-4$  & 1 & trivial   & 1 & 1 (G2 ssh8)    \\
$-7$  & 1 & trivial   & 1 & 1 (G2 ssh8)    \\
$-15$ & 2 & $\Z/2$    & 2 & 2 (Co8v2)      \\
$-20$ & 2 & $\Z/2$    & 2 & 2 (Co8v2)      \\
$-23$ & 3 & $\Z/3$    & 1 & 1 canonical only (Bv6) \\
$-31$ & 3 & $\Z/3$    & 1 & 1 canonical only (Bv6) \\
$-59$ & 3 & $\Z/3$    & 1 & 1 canonical only (Bv6) \\
$-83$ & 3 & $\Z/3$    & 1 & 1 canonical only (Bv6) \\
$-107$& 3 & $\Z/3$    & 1 & 1 canonical only (Bv6) \\
$-39$ & 4 & $\Z/4$    & 2 & 2 (Bv3 partial) \\
$-84$ & 4 & $(\Z/2)^2$ & 3 & \emph{Bv7 target} \\
$-420$& 8 & $(\Z/2)^3$ & 4 & \emph{Bv7 target} \\
\bottomrule
\end{tabular}
\caption{Numerical evidence for Conjecture~\ref{conj:HSH} at weight $w=5$.
The $h_K \in \{1,2,3\}$ rows are confirmed; the $\rktwo = 2$ biquadratic row and the
$\rktwo = 3$ row are open and constitute the principal falsifier (Bv7 sweep,
Section~\ref{ss:falsifier}).}
\label{tab:Bv6-evidence}
\end{table}

The five $h_K = 3$ rows are the empirical seed of the conjecture: each was a candidate
for two non-trivial $\Q$-rational newforms (one for $\chi$ and one for $\chi^2$, or one
$\Q$-rational trace), and the actual observation was \emph{zero} extra newforms beyond
the canonical theta-series. The match with the prediction $\rktwo + 1 = 1$ is the
$8/8$ figure of the abstract.

\subsection{The Bv7 falsifier}\label{ss:falsifier}

The principal open test of Conjecture~\ref{conj:HSH} is the biquadratic case
$\Cl(K) = (\Z/2)^2$, in particular $D \in \{-84, -120, -132, -168, -228\}$. The HSH
prediction is $\rktwo + 1 = 3$. We refer to the sweep that tests these anchors as Bv7;
the sweep is CPU-bound (PARI \texttt{mfinit} at weight $5$, level $\sim 10^4$) and is
expected to take roughly $20$--$45$ minutes on a 96-core EPYC instance with parallelised
discriminants. We also include $\rktwo = 3$ anchors ($D \in \{-420, -660, -780\}$)
predicting $\dim = 4$, and $\rktwo = 1, h_K > 2$ controls ($D \in \{-39, -55, -56\}$)
predicting $\dim = 2$.

The conjecture is \emph{falsified} if any biquadratic anchor gives $\dim \le 2$, or if
any $\rktwo = 0, h_K$-odd, $h_K > 1$ anchor gives $\dim \ge 2$. A weaker degradation
to $\dim = \min(\rktwo, 1) + 1$ would preserve the $h_K \le 2$ dichotomy of the title
but lose the higher-rank structure.

\subsection{Derivation sketch and the residual lemma}\label{ss:derivation}

The derivation of \eqref{eq:HSH} proceeds in five steps; details would be filled out in
a full version of this paper.

\begin{enumerate}[label=(S\arabic*),leftmargin=*]
\item By~\cite{Schutt2009}, $\Q$-rationality of $f_\psi$ forces $\psi$ to factor through
the genus quotient $\Cl(K)/\Cl(K)^2 \cong (\Z/2)^{\rktwo}$.
\item The character group of $\Cl(K)/\Cl(K)^2$ has cardinality $2^{\rktwo}$
(Pontryagin duality).
\item Two characters give twist-equivalent newforms over $\Q$ iff their ratio is a
Dirichlet character of $\Q$; by Hecke~\cite{Hecke1936} (revisited via the
ring-class-field formalism of Schertz~\cite{Schertz1976}), the $\Q$-Dirichlet quotient
of the genus character group has index $\rktwo + 1$.
\item Hence the number of $\Q$-twist classes is $\rktwo + 1$.
\item The Schertz-sparse level union $\Lcal(K, f_{\max})$ with $f_{\max} \ge 6$ detects
all $\Q$-twist classes (Schertz visibility; empirically verified up to $\hK = 4$).
\end{enumerate}

Steps (S1)--(S3) and (S5) are rigorous. Step (S4) elevates a count-quotient
identification to a \emph{Lemma~A}, the explicit identification of the
$\Q$-Dirichlet-twist quotient. The lemma is "morally clear" but its tightening into a
formal proof is the principal theoretical task remaining; we flag it as the gating step
of any JHEP-or-CMP final submission.

\section{Connections}\label{sec:connections}

\subsection{Borcea--Voisin orbifolds}\label{ss:Borcea-Voisin}

The Borcea--Voisin construction~\cite{Borcea1997, Voisin1993} produces a Calabi--Yau
threefold from a K3 surface $X$ with non-symplectic involution $\sigma$ and an elliptic
curve $E$ with $-\mathrm{id}$:
\[
  Y_{BV} \;=\; \widetilde{(X \times E)/(\sigma \times -\mathrm{id})}.
\]
The construction works because $\sigma$ and $-\mathrm{id}$ are \emph{involutions} (order 2):
the holomorphic 3-form $\omega_X \wedge \omega_E$ is sent to itself.

The analogous order-3 construction $(X \times E)/(\sigma_3 \times \rho_3)$ generally
fails to give a Calabi--Yau, since $\sigma_3$ acts on $H^{0,2}(X)$ by $\zeta_3$, not by
$-1$. This geometric obstruction is the structural analogue of the arithmetic
obstruction in Conjecture~\ref{conj:HSH}: order-$2$ quotients preserve $\Q$-rational
structure, order-$3$ quotients do not. The parallel is striking and we conjecture (with
lower confidence) that the Borcea--Voisin construction lifts to a $K3$-fibred bulk in
$\AdS_3$ whose boundary $\CFT_2$ encodes precisely the $\rktwo + 1$ count of
Conjecture~\ref{conj:HSH}.

\subsection{Birch--Tate and Iwasawa BCS dichotomy}\label{ss:BirchTate}

Conjecture~\ref{conj:HSH} sits naturally alongside an Iwasawa-side companion. The
Birch--Tate conjecture predicts $|K_2(\Ocal_F)|$ in terms of $\zeta_F(-1)$ for totally
real $F$, with the order of the wild-kernel subgroup encoded by class group $2$-torsion.
For CM elliptic curves $E/\Q$ with CM by $K$, the Selmer group $\mathrm{Sel}_p(K_\infty/K)$
exhibits a phase transition controlled by the parity of $\hK$: $\hK$ odd implies a
trivial $2$-torsion in the Selmer characteristic ideal, an "Iwasawa BCS superconducting"
phase. Conjecture~\ref{conj:HSH} is the \emph{modular-side mirror} of this Selmer-side
dichotomy: in both settings, the obstruction is parametrised by $\rktwo(\Cl(K))$.

\subsection{Connes--Marcolli IC8 fatgraph}\label{ss:IC8}

The Connes--Marcolli BC type III$_1$ programme~\cite{ConnesMarcolli2008} attaches a
noncommutative-geometric structure to the modular flow of class field theory. Internal
ECI v12 morn39 day-end analysis identifies an inner-component-$8$ (IC8) fatgraph with
central charge $c = 4/5$ as a natural NCG model for the $\Z/3$-orbifold boundary CFT.
The coincidence $c = 4/5 =$ three-state Potts central charge is, we propose, structural:
both encode the $\Z/3$ obstruction discussed in Section~\ref{ss:Zk-obstruction}, in
two different guises. A precise IC8-Potts equivalence would constitute an HSH-NCG
bridge, parallel to and reinforcing the HSH conjecture itself.

\subsection{Tamagawa cross-check}\label{ss:Tamagawa}

Internal ECI v12 morn39 day-end analysis (Tamagawa = $N_W$ bridge, $6/6$ exact) suggests
a second cross-check for Conjecture~\ref{conj:HSH}: for CM elliptic curves $E/\Q$ with
CM by $\OK$, the Tamagawa-side dimension of the bad-reduction-supported modular subspace
should match $\rktwo + 1$. If Bv7 confirms the biquadratic prediction, the Tamagawa
identity should hold for all anchors with the corresponding multiplicity, giving an
independent verification path that does not invoke the holographic dictionary directly.
We flag this Tamagawa-side cross-check as a separate target for Section~\ref{ss:next-steps}.

\section{Open problems and outlook}\label{sec:open-problems}

\subsection{Anticipated objections}\label{ss:objections}

\textbf{Obj.\ 1 (string-theory reviewer):} \emph{the holographic connection is structural
analogy, not duality.} Indeed; we frame HSH as a structural orbifold-projection count,
not a derived bulk--boundary duality. A quantitative match would require either an exact
bulk Lagrangian whose path integral computes $\rktwo + 1$, or a categorical lift via the
Connes--Marcolli BC programme.

\textbf{Obj.\ 2 (analytic number-theory reviewer):} \emph{Sch\"utt's classification is
unconditional only through weight 4.} The Bv6 sweep at weight 5 is purely empirical; the
weight-5 HSH is a prediction, tested by~Table~\ref{tab:Bv6-evidence} and to be tested by
Bv7--Bv10 (see Section~\ref{ss:next-steps}).

\textbf{Obj.\ 3 (skeptical Langlands person):} \emph{this looks like automorphic-induction
counting.} Yes; HSH is precisely the assertion that the automorphic-induction count at
the Schertz-sparse level union equals $\rktwo + 1$, plus the additional structural
claim that this count is the holographic shadow of a $\Z/2$ orbifold projection.

\subsection{Open theoretical questions}\label{ss:openQ}

\begin{enumerate}[label=Q\arabic*.,leftmargin=*]
\item \textbf{Cohomological formula.} Is there a cohomological count
$\dim H^*(\mathrm{some\ local\ system\ on\ }\Cl(K)) = \rktwo + 1$ giving an \'etale or
crystalline derivation of Conjecture~\ref{conj:HSH}?
\item \textbf{Real-quadratic extension.} Does an HSH-RM analogue hold for $K = \Q(\sqrt{D})$,
$D > 0$, using Hilbert modular forms of parallel weight $w$?
\item \textbf{Discriminant dependence beyond $\rktwo$.} Does the formula see $h_K$ at all,
or only $\rktwo$? Empirically $h_K = 3$ and $h_K = 5$ both predict $\dim = 1$;
Bv8 anchors $D \in \{-47, -79, -103, -127\}$ ($h_K = 5$) and $D \in \{-71, -151, -223\}$
($h_K = 7$) will test this.
\item \textbf{Heegner-point subspace.} Do the rational CM newforms detected by HSH
correspond to the modular forms whose Heegner points have $\Q$-rational coordinates?
The Gross--Zagier framework suggests yes.
\item \textbf{Motivic upgrade.} Is HSH the decategorified shadow of a Tannakian statement
in Deligne's category of CM motives over $\Q$?
\end{enumerate}

\subsection{Risk register}\label{ss:risks}

We collect six principal risks of the present draft.

\begin{itemize}
\item[R1.] Bv7 biquadratic gives $\dim = 2$ (not 3). Mitigation: try $f_{\max} \ge 8$
(probability $20\%$).
\item[R2.] $\Q$-twist equivalence is mis-identified; actual count is $2^{\rktwo}$ rather
than $\rktwo + 1$. Mitigation: explicit twist-character enumeration ($15\%$).
\item[R3.] The $\AdS_3/\CFT_2$ dictionary is metaphor rather than quantitative match.
Mitigation: literature search for a precise prior-art bridge; reframe as analogy
($35\%$).
\item[R4.] IC8 $c = 4/5$ match is coincidence. Mitigation: explicit Connes--Marcolli
IC8 vs.\ Potts comparison ($25\%$).
\item[R5.] Sch\"utt 2009 already implies a stronger statement. Mitigation: cite~\cite{Schutt2009}
explicitly; HSH refines finiteness to an explicit count ($10\%$).
\item[R6.] Reviewer skepticism on the "Holographic Schutt-Hecke" branding. Mitigation:
reframe as "an arithmetic refinement of Sch\"utt's finiteness conjecture, with
conjectural $\AdS_3/\CFT_2$ counterparts" ($30\%$).
\end{itemize}

\subsection{Next steps}\label{ss:next-steps}

Three concrete next actions are scheduled for completion within ECI v15.

\begin{enumerate}[label=(\arabic*),leftmargin=*]
\item \textbf{Bv7 PARI sweep} on a 96-core EPYC instance, $\rktwo \in \{2,3\}$ anchors
and $\rktwo = 1, h_K > 2$ controls; expected runtime $20$--$45$ min.
\item \textbf{Literature search} on prior bridges between CM modular forms and
$\AdS_3/\CFT_2$ / chiral $\CFT_2$ boundary conditions; targets include Manin's
noncommutative-geometry-and-number-theory programme and the Connes--Marcolli $\Q$-lattice
formalism.
\item \textbf{Full paper expansion.} Conditional on Bv7 confirmation, expansion of the
present outline to a 25--35 page submission for JHEP (primary) or Communications in
Mathematical Physics (secondary).
\end{enumerate}

\subsection{Outlook}\label{ss:outlook}

If Conjecture~\ref{conj:HSH} survives the Bv7 falsifier, the natural next direction is
twofold. First, one can ask whether the $\rktwo + 1$ count admits a categorical
upgrade in the spirit of Deligne's category of motives over $\Q$ with CM by $K$,
producing a Tannakian statement whose dimension count is \eqref{eq:HSH}. Second, one
can ask whether the holographic dictionary of Section~\ref{sec:AdS3-CFT2} admits a
quantitative match at the level of partition functions: a precise bulk Lagrangian
whose partition function on $\AdS_3 \times M^4$, with boundary conditions selecting
the Schertz-sparse level union, reproduces the count $\rktwo + 1$ along with the
Tamagawa cross-check of Section~\ref{ss:Tamagawa}. Both directions strengthen the
case for HSH as a genuine arithmetic-holographic principle rather than a structural
analogy.

A third, more speculative direction concerns the \emph{real} quadratic case
$K = \Q(\sqrt D)$ with $D > 0$, where CM is replaced by real multiplication and the
theta-series of Section~\ref{sec:Schutt-Hecke} are replaced by Hilbert modular forms
of parallel weight $w$. Genus theory still applies, and the analogue of $\rktwo$ is
again read off from the prime factorisation of $d_K$. We conjecture that an HSH-RM
analogue holds, with the same $\rktwo + 1$ count, though the appropriate analogue of the
Schertz-sparse level union and the appropriate notion of $\iota$-invariant character
need to be defined carefully.

\appendix

\section{PARI/GP pseudocode for Bv7}\label{app:Bv7-pseudo}

For reproducibility, we record a sketch of the Bv7 sweep used to test
Conjecture~\ref{conj:HSH} at $\rktwo \in \{1, 2, 3\}$. The actual sweep uses additional
caching and parallelisation; the core logic is:

\begin{verbatim}
read("Bv6_opt_core.gp");
results = Mat();
for(D in [-84,-120,-132,-168,-228, -420,-660,-780, -39,-55,-56],
  K = bnfinit(x^2 + D);
  hK = K.cyc;
  rk2 = #select(c -> c%2==0, hK);
  L = sparse_level_union(D, fmax=6);
  rats = 0;
  for(N in L,
    S = mfinit([N, 5], 0);
    for(f in mfbasis(S),
      if(is_CM(f, D) && is_Q_rational(f)
         && !already_twist_seen(f, results),
         rats += 1;
         results = matadjoin(results, [D, N, f])
      )
    );
  );
  print("D=", D, " rk2=", rk2,
        " predicted=", rk2+1, " observed=", rats);
);
\end{verbatim}

The expected output, under Conjecture~\ref{conj:HSH}, is "observed = predicted" on every
line. A discrepancy at $D \in \{-84, -120, -132, -168, -228\}$ would falsify the
biquadratic case; a discrepancy at $D \in \{-420, -660, -780\}$ would falsify the
$\rktwo = 3$ case.

\section*{Acknowledgements}

This work emerged from the multi-LLM ECI v12--v14 research programme (Hostinger VPS,
April--May 2026). The author thanks the various LLM dispatch threads for collaborative dispatches under explicit
anti-fabrication discipline; all numerical, bibliographic, and identifier claims in this
draft have been independently verified.

\begin{thebibliography}{99}

\bibitem{BTZ1992}
M.\ Ba\~nados, C.\ Teitelboim, J.\ Zanelli,
"The Black Hole in Three Dimensional Space Time",
\emph{Phys.\ Rev.\ Lett.} \textbf{69} (1992), 1849--1851,
arXiv:hep-th/9204099.

\bibitem{Borcea1997}
C.\ Borcea,
"K3 surfaces with involution and mirror pairs of Calabi-Yau manifolds",
\emph{Mirror Symmetry II}, AMS/IP Stud.\ Adv.\ Math.\ \textbf{1} (1997), 717--743.

\bibitem{BrownHenneaux1986}
J.\ D.\ Brown, M.\ Henneaux,
"Central charges in the canonical realization of asymptotic symmetries: An example
from three dimensional gravity",
\emph{Commun.\ Math.\ Phys.} \textbf{104} (1986), 207--226.

\bibitem{ConnesMarcolli2008}
A.\ Connes, M.\ Marcolli,
\emph{Noncommutative Geometry, Quantum Fields and Motives},
AMS Colloquium Publications \textbf{55}, 2008.

\bibitem{Deligne1979}
P.\ Deligne,
"Valeurs de fonctions $L$ et p\'eriodes d'int\'egrales",
\emph{Proc.\ Symp.\ Pure Math.} \textbf{33} (1979), 313--346.

\bibitem{Gauss1801}
C.\ F.\ Gauss,
\emph{Disquisitiones Arithmeticae},
\S 305 (genus theory), 1801.

\bibitem{HarrisKudla1991}
M.\ Harris, S.\ Kudla,
"The central critical value of a triple product $L$-function",
\emph{Ann.\ of Math.} \textbf{133} (1991), 605--672.

\bibitem{Hecke1936}
E.\ Hecke,
\emph{Mathematische Werke}, \S 2, 1936.

\bibitem{Maldacena1997}
J.\ M.\ Maldacena,
"The Large $N$ Limit of Superconformal Field Theories and Supergravity",
\emph{Adv.\ Theor.\ Math.\ Phys.} \textbf{2} (1998), 231--252,
arXiv:hep-th/9711200.

\bibitem{MaldacenaStrominger1998}
J.\ M.\ Maldacena, A.\ Strominger,
"$\AdS_3$ Black Holes and a Stringy Exclusion Principle",
\emph{JHEP} \textbf{12} (1998), 005,
arXiv:hep-th/9804085.

\bibitem{MooreSeiberg1989}
G.\ Moore, N.\ Seiberg,
"Classical and quantum conformal field theory",
\emph{Commun.\ Math.\ Phys.} \textbf{123} (1989), 177--254.

\bibitem{Schertz1976}
R.\ Schertz,
"Niedere Potenzen elliptischer Funktionen und Klassenk\"orpertheorie",
\emph{J.\ Reine Angew.\ Math.} \textbf{285} (1976), 49--65.

\bibitem{Schutt2009}
M.\ Sch\"utt,
"CM newforms with rational coefficients",
\emph{Ramanujan J.} \textbf{19} (2009), 187--205,
arXiv:math/0511228.

\bibitem{Strominger1998}
A.\ Strominger,
"Black Hole Entropy from Near-Horizon Microstates",
\emph{JHEP} \textbf{02} (1998), 009,
arXiv:hep-th/9712251.

\bibitem{Verlinde1988}
E.\ Verlinde,
"Fusion rules and modular transformations in 2D conformal field theory",
\emph{Nucl.\ Phys.\ B} \textbf{300} (1988), 360--376.

\bibitem{Voisin1993}
C.\ Voisin,
"Miroirs et involutions sur les surfaces K3",
in \emph{Journ\'ees de G\'eom\'etrie Alg\'ebrique d'Orsay}, Ast\'erisque \textbf{218}
(1993), 273--323.

\bibitem{Witten1989}
E.\ Witten,
"Quantum field theory and the Jones polynomial",
\emph{Commun.\ Math.\ Phys.} \textbf{121} (1989), 351--399.

\end{thebibliography}

\end{document}
